\clearpage
\oldpage[IV-2]{229}
\phantomsection
\addcontentsline{toc}{section}{Tabla de contenido}
\section*{TABLA DE CONTENIDO}

\begingroup
\small
\setlength{\parindent}{0pt}
\setlength{\parskip}{1.5pt}
\hfill\textsc{Páginas}\par
\newcommand{\EGAIVtwoContentsLine}[3]{#1\quad #2\dotfill #3\par}

\EGAIVtwoContentsLine{\textsc{Capítulo IV.}}{--- \textbf{Estudio local de los esquemas y de los
morfismos de esquemas} \emph{(continuación)}}{5}

\EGAIVtwoContentsLine{\S~2.}{Cambio de base y planitud}{5}
\EGAIVtwoContentsLine{2.1.}{Módulos planos sobre los preesquemas}{5}
\EGAIVtwoContentsLine{2.2.}{Módulos fielmente planos sobre los preesquemas}{9}
\EGAIVtwoContentsLine{2.3.}{Propiedades topológicas de los morfismos planos}{14}
\EGAIVtwoContentsLine{2.4.}{Morfismos universalmente abiertos y morfismos planos}{19}
\EGAIVtwoContentsLine{2.5.}{Permanencia de las propiedades de los módulos por descenso fielmente plano}{22}
\EGAIVtwoContentsLine{2.6.}{Permanencia de propiedades conjuntistas y topológicas de morfismos por descenso fielmente plano}{27}
\EGAIVtwoContentsLine{2.7.}{Permanencia de diversas propiedades de los morfismos por descenso fielmente plano}{29}
\EGAIVtwoContentsLine{2.8.}{Preesquemas sobre una base regular de dimensión 1; clausura de un subpreesquema cerrado de la fibra genérica}{33}

\EGAIVtwoContentsLine{\S~3.}{Ciclos primos asociados y descomposiciones primarias}{36}
\EGAIVtwoContentsLine{3.1.}{Ciclos primos asociados a un módulo}{36}
\EGAIVtwoContentsLine{3.2.}{Descomposiciones irredundantes}{40}
\EGAIVtwoContentsLine{3.3.}{Relaciones con la planitud}{43}
\EGAIVtwoContentsLine{3.4.}{Propiedades de los haces $\mathscr F/t\mathscr F$}{46}

\EGAIVtwoContentsLine{\S~4.}{Cambio del cuerpo base en los preesquemas algebraicos}{52}
\EGAIVtwoContentsLine{4.1.}{Dimensión de los preesquemas algebraicos}{52}
\EGAIVtwoContentsLine{4.2.}{Ciclos primos asociados sobre los preesquemas algebraicos}{54}
\EGAIVtwoContentsLine{4.3.}{Recordatorio sobre los productos tensoriales de cuerpos}{58}
\EGAIVtwoContentsLine{4.4.}{Preesquemas irreducibles y preesquemas conexos sobre un cuerpo algebraicamente cerrado}{59}
\EGAIVtwoContentsLine{4.5.}{Preesquemas geométricamente irreducibles y geométricamente conexos}{61}
\EGAIVtwoContentsLine{4.6.}{Preesquemas algebraicos geométricamente reducidos}{68}

\clearpage
\oldpage[IV-2]{230}

\hfill\textsc{Páginas}

\EGAIVtwoContentsLine{4.7.}{Multiplicidades en la descomposición primaria sobre un preesquema algebraico}{75}
\EGAIVtwoContentsLine{4.8.}{Cuerpos de definición}{80}
\EGAIVtwoContentsLine{4.9.}{Cuerpos de definición de un subconjunto de un preesquema}{84}

\EGAIVtwoContentsLine{\S~5.}{Dimensión, profundidad y regularidad en los preesquemas localmente noetherianos}{86}
\EGAIVtwoContentsLine{5.1.}{Dimensión de los preesquemas}{86}
\EGAIVtwoContentsLine{5.2.}{Dimensión de un preesquema algebraico}{90}
\EGAIVtwoContentsLine{5.3.}{Dimensión del soporte de un módulo y polinomio de Hilbert}{92}
\EGAIVtwoContentsLine{5.4.}{Dimensión de la imagen de un morfismo}{93}
\EGAIVtwoContentsLine{5.5.}{Fórmula de las dimensiones para un morfismo de tipo finito}{94}
\EGAIVtwoContentsLine{5.6.}{Fórmula de las dimensiones y anillos universalmente catenarios}{97}
\EGAIVtwoContentsLine{5.7.}{Profundidad y propiedad $(S_k)$}{103}
\EGAIVtwoContentsLine{5.8.}{Preesquemas regulares y propiedad $(R_k)$. Criterio de normalidad de Serre}{107}
\EGAIVtwoContentsLine{5.9.}{Módulos $Z$-puros y $Z$-cerrados}{109}
\EGAIVtwoContentsLine{5.10.}{Propiedad $(S_2)$ y $Z$-clausura}{114}
\EGAIVtwoContentsLine{5.11.}{Criterios de coherencia para los módulos $\mathscr H^0_{X/Z}(\mathscr F)$}{122}
\EGAIVtwoContentsLine{5.12.}{Relaciones entre las propiedades de un anillo local noetheriano $A$ y de un anillo cociente $A/tA$}{126}
\EGAIVtwoContentsLine{5.13.}{Propiedades de permanencia al pasar al límite inductivo}{131}

\EGAIVtwoContentsLine{\S~6.}{Morfismos planos de preesquemas localmente noetherianos}{134}
\EGAIVtwoContentsLine{6.1.}{Planitud y dimensión}{135}
\EGAIVtwoContentsLine{6.2.}{Planitud y dimensión proyectiva}{137}
\EGAIVtwoContentsLine{6.3.}{Planitud y profundidad}{138}
\EGAIVtwoContentsLine{6.4.}{Planitud y propiedad $(S_k)$}{141}
\EGAIVtwoContentsLine{6.5.}{Planitud y propiedad $(R_k)$}{143}
\EGAIVtwoContentsLine{6.6.}{Propiedades de transitividad}{145}
\EGAIVtwoContentsLine{6.7.}{Aplicación a los cambios de base en los preesquemas algebraicos}{145}
\EGAIVtwoContentsLine{6.8.}{Morfismos regulares, normales, reducidos y suaves}{150}
\EGAIVtwoContentsLine{6.9.}{El teorema de planitud genérica}{153}
\EGAIVtwoContentsLine{6.10.}{Dimensión y profundidad de un módulo normalmente plano a lo largo de un subpreesquema cerrado}{155}
\EGAIVtwoContentsLine{6.11.}{Criterios para que los conjuntos $U_{S_n}(\mathscr F)$ o $U_{C_n}(\mathscr F)$ sean abiertos}{158}
\EGAIVtwoContentsLine{6.12.}{Criterios de Nagata para que $\operatorname{Reg}(X)$ sea abierto}{163}
\EGAIVtwoContentsLine{6.13.}{Criterios para que $\operatorname{Nor}(X)$ sea abierto}{168}
\EGAIVtwoContentsLine{6.14.}{Cambio de base y clausura entera}{169}
\EGAIVtwoContentsLine{6.15.}{Preesquemas geométricamente unirrama}{176}

\clearpage
\oldpage[IV-2]{231}

\hfill\textsc{Páginas}

\EGAIVtwoContentsLine{\S~7.}{Relaciones entre un anillo local noetheriano y su completado. Anillos excelentes}{182}
\EGAIVtwoContentsLine{7.1.}{Equidimensionalidad formal y anillos formalmente catenarios}{183}
\EGAIVtwoContentsLine{7.2.}{Anillos estrictamente formalmente catenarios}{187}
\EGAIVtwoContentsLine{7.3.}{Fibras formales de los anillos locales noetherianos}{192}
\EGAIVtwoContentsLine{7.4.}{Permanencia de las propiedades de las fibras formales}{198}
\EGAIVtwoContentsLine{7.5.}{Un criterio para los $\boldsymbol P$-morfismos}{203}
\EGAIVtwoContentsLine{7.6.}{Aplicaciones: I. Anillos japoneses locales}{208}
\EGAIVtwoContentsLine{7.7.}{Aplicaciones: II. Anillos universalmente japoneses}{212}
\EGAIVtwoContentsLine{7.8.}{Anillos excelentes}{214}
\EGAIVtwoContentsLine{7.9.}{Anillos excelentes y resolución de singularidades}{218}
\EGAIVtwoContentsLine{}{\textsc{Bibliografía} \emph{(continuación)}}{224}
\EGAIVtwoContentsLine{}{\textsc{Índice de notaciones}}{225}
\EGAIVtwoContentsLine{}{\textsc{Índice terminológico}}{226}

\vspace{2em}
\begin{center}
\emph{Manuscrito recibido el 15 de diciembre de 1963.}
\end{center}

\begin{center}\rule{25mm}{0.4pt}\end{center}

\endgroup

\par\smallskip
\noindent\textit{Nota de lectura.} Las páginas de esta tabla corresponden al
impreso original, no a la paginación de los lectores traducidos.
